%% kapa: exam=1 round=32 year=2021 period=1 form=A subject=civil_law subject_ko=민법 seq=32 number=32
%% source: Q-Net 감정평가사 기출문제 (cst00309) artlSeq=5208397 file=extracted/1008/5208397/A형/2021년도 제32회 감정평가사 1차 1교시 A형.pdf
%% conversion: kordoc 4.24.51 -> tools/gen1.py (auto); review=proofread-sonnet; flags=none
\begin{question}{32}
甲은 乙에 대한 채권을 담보하기 위하여 乙 소유의 X토지에 관하여 저당권을 취득하였다. 그 후 X의 담보가치 하락을 막기 위하여 乙의 X에 대한 사용·수익권을 배제하지 않는 지상권을 함께 취득하였다. 이에 관한 설명으로 옳지 않은 것은? (다툼이 있으면 판례에 따름)
\begin{choices}
  \item 甲의 지상권의 피담보채무는 존재하지 않는다.
  \item 甲의 채권이 시효로 소멸하면 지상권도 소멸한다.
  \item 甲의 채권이 변제 등으로 만족을 얻어 소멸하면 지상권도 소멸한다.
  \item 제3자가 甲에게 대항할 수 있는 권원 없이 X 위에 건물을 신축하는 경우, 甲은 그 축조의 중지를 요구할 수 있다.
  \item 제3자가 X를 점유·사용하는 경우, 甲은 지상권의 침해를 이유로 손해배상을 청구할 수 있다.
\end{choices}
\end{question}
\answer{5}
